% Source: https://artofproblemsolving.com/community/c3h2737679_3rd_and_last_3500th_post_celebration_
% Stable identity: aops-topic-2737679-q01
Since $D$ lies on line $AB$, $\angle DBC=90^\circ$. An inscribed right angle subtends a diameter, so $DC$ is a diameter and $O$ is its midpoint.

\begin{asy}
usepackage("amsmath");
size(4.5cm);
defaultpen(fontsize(10pt));
real t=40;
pair O=(0,0),C=(1,0),D=(-1,0),B=dir(130);
pair A=D+2*(D-B);
pair[] X=intersectionpoints(A--C,unitcircle);
pair E=(abs(X[0]-C)>1e-6)?X[0]:X[1];
pair M=(D+E)/2;
draw(unitcircle,gray);draw(A--B--C--cycle,linewidth(0.8));draw(D--C,dashed);draw(D--E);draw(O--M,linewidth(1));
dot(O);dot(M);dot(D);dot(E);
label("$A$",A,SW);label("$B$",B,N);label("$C$",C,E);label("$D$",D,W);label("$E$",E,S);label("$O$",O,N);label("$M$",M,NE);
\end{asy}

In triangle $DCE$, $O$ is the midpoint of $DC$ and $M$ is the midpoint of $DE$, so $OM$ is a midline and $OM\parallel CE$. Since $E$ lies on $AC$, line $CE$ is line $AC$, so $OM\parallel AC$. $\blacksquare$
